\section{Discussion and Conclusion}
\label{sec:conclusion}
%
\paragraph{Limitations and future work}
Our stochastic formulation introduces variance into the gradient estimates, which we observed affects the behavior of the splitting and pruning heuristics used for Gaussian densification. A more detailed investigation of how the densification scheme interacts with this variance is left to future work.

\paragraph{Conclusion}
We presented a differentiable, sorting-free stochastic ray tracing framework for 3D Gaussian Splatting---the first to use stochastic ray tracing for both reconstruction and rendering of standard and relightable scenes.
By replacing exhaustive sorted evaluation with an unbiased Monte Carlo estimator that samples only a small subset of Gaussians per ray, our method matches the speed and quality of rasterization-based 3DGS while substantially outperforming sorting-based ray tracing.
For relightable scenes, the same estimator enables per-Gaussian shading with fully ray-traced shadow rays, delivering notably higher reconstruction fidelity than prior approaches that rely on rasterization-style approximations.

\section*{Acknowledgments}
We thank the anonymous reviewers for their feedback and suggestions.
This work was partially supported by NSF grant 2553564. This work started when Peiyu Xu was an intern at Adobe Research.
